\documentclass[10pt]{article}
\usepackage[margin=1.35cm]{geometry}
\usepackage{booktabs,amsmath,array,tabularx}
\setlength{\emergencystretch}{2em}
\title{Full-Signal PNNN versus PN-GMP DOMP Comparison}
\author{}\date{}\begin{document}\maketitle
\paragraph{Objective and methodology.} Complex GMP, its exactly equivalent coupled PN-IQ representation, independent PN-IQ controls, and parameter-matched PNNNs use the same modeled-block X/Y samples. DOMP is used only during linear identification.
\paragraph{Fair comparison protocol.} The classic selector provides 49151 identification rows (configured 10\% of the capture). An internal 41778/7373 split selects regularization and neural epochs. Every final model is fitted on all identification rows and evaluated on all 491520 signal rows. \textbf{The full-signal evaluation includes the identification samples and therefore must not be interpreted as an independent generalization test.}
\paragraph{Main results.} NMSE uses $10\log_{10}(\sum|y-\hat y|^2/\sum|y|^2)$.
\begin{center}\scriptsize\begin{tabularx}{\linewidth}{XllrrrX}
\toprule
Model & Regularization & $\lambda$ & Parameters & Full-signal NMSE (dB) & FLOPs/sample & Additional operations \\ \midrule
Independent PN-IQ full & LS & 0 & 344 & -38.839 & 1026 & 13 magnitude (including 13 sqrt), 2 division, phase normalization/restoration arithmetic in FLOPs \\ 
Independent PN-IQ full & Fixed Ridge & 1e-05 & 344 & -37.120 & 1026 & 13 magnitude (including 13 sqrt), 2 division, phase normalization/restoration arithmetic in FLOPs \\ 
Independent PN-IQ full & Validation-selected Ridge & 0 & 344 & -38.839 & 1026 & 13 magnitude (including 13 sqrt), 2 division, phase normalization/restoration arithmetic in FLOPs \\ 
Complex GMP DOMP parameter-matched & LS & 0 & 344 & -38.458 & 1837 & 14 magnitude (including 14 sqrt) \\ 
Complex GMP DOMP parameter-matched & Fixed Ridge & 1e-05 & 344 & -36.860 & 1837 & 14 magnitude (including 14 sqrt) \\ 
Complex GMP DOMP parameter-matched & Validation-selected Ridge & 0 & 344 & -38.458 & 1837 & 14 magnitude (including 14 sqrt) \\ 
PNNN N12 sparse & Not applicable & -- & 344 & -35.660 & 848 & 14 magnitude (including 14 sqrt), 2 division, 12 sigmoid, up to 12 exp, phase normalization/restoration arithmetic in FLOPs \\ 
\bottomrule\end{tabularx}\end{center}
\paragraph{Computational cost.} FLOPs/sample counts real additions and multiplications required for one complex output, including feature generation, coefficients or weights, biases, phase normalization, and phase restoration. Magnitudes, roots, divisions, exponentials, and activation evaluations remain separate. DOMP is not an inference cost. The sparse-PNNN FLOP count assumes that zero weights are skipped. An implementation that executes the original full matrices will require 2252 FLOPs/sample instead.
\paragraph{Interpretation and limitations.} Independent PN-IQ provides the strongest NMSE in the main comparison; all three main models use exactly the same number of active real parameters; and the sparse PNNN uses the fewest counted FLOPs when zero weights are actually skipped, while adding sigmoid evaluations. Complex GMP DOMP-100 and coupled PN-IQ remain numerically equivalent for LS, fixed Ridge, and validation-selected Ridge, with full-signal relative errors $1.326e-11$, $7.661e-16$, and $1.326e-11$. These analytical counts do not establish latency, energy, or FPGA resource superiority.
\end{document}
